\documentclass[a4paper,11pt]{article}
\usepackage[utf8]{inputenc}
\usepackage[french]{babel}
\usepackage[T1]{fontenc}
\usepackage{amsmath,amssymb,amsthm}
\usepackage{geometry}
\geometry{margin=2cm}
\usepackage{enumitem}
\usepackage{hyperref}
\usepackage{booktabs}
\usepackage{xcolor}
\usepackage{listings}
\usepackage{titlesec}
\usepackage{fancyhdr}
\usepackage{lastpage}
\AtBeginDocument{\shorthandoff{;:!?}}

\titleformat{\section}{\large\bfseries\color{blue!60!black}}{}{0pt}{}[\vspace{-0.5ex}\rule{\textwidth}{0.4pt}]
\titleformat{\subsection}{\bfseries\color{blue!40!black}}{}{0pt}{}
\lstdefinestyle{pystyle}{
  frame=single, numbers=left, breaklines=true,
  basicstyle=\small\ttfamily, backgroundcolor=\color{gray!5},
  rulecolor=\color{gray!40}, language=Python
}
\lstset{style=pystyle}

\pagestyle{fancy}
\fancyhf{}
\fancyhead[L]{\small STA211 -- Notebook \'enonc\'e (Keras)}
\fancyhead[R]{\small Fiche r\'ecapitulative}
\fancyfoot[C]{\thepage/\pageref{LastPage}}
\renewcommand{\headrulewidth}{0.4pt}

\title{\textbf{Fiche r\'ecapitulative -- STA211}\\[4pt]
\large Notebook Jupyter -- \'Enonc\'e \\[2pt]
\small DL-STA211.ipynb -- Deep Learning avec Keras}
\author{N. Thome (CNAM)}
\date{14 juin 2026}

\begin{document}
\maketitle
\thispagestyle{empty}

\begin{abstract}
Le notebook \'enonc\'e fournit aux \'etudiants le squelette du TP
Keras/Manifold Untangling avec des cellules de code \`a compl\'eter
et des questions \`a r\'epondre. Cette fiche r\'esume la structure,
les cellules de code et les questions du notebook.
\end{abstract}

\vfill
\tableofcontents
\newpage

% ===================================================================
\section{Structure du notebook}
% ===================================================================

Le notebook contient 67 cellules (markdown + code) couvrant 6
exercices. Les cellules de code ont \texttt{execution\_count=None}
(doivent \^etre ex\'ecut\'ees par l'\'etudiant).

\subsection{Partie 1 -- Deep Learning avec Keras}

\paragraph{Ex0 -- Chargement des donn\'ees (2 cellules code)~:}
charger MNIST, normaliser, reshape en 784, afficher 200 images.

\paragraph{Ex1 -- R\'egression logistique (5 cellules code)~:}
\begin{enumerate}
  \item Cr\'eer un \texttt{Sequential} vide
  \item Ajouter \texttt{Dense(10, input\_dim=784)} + softmax
  \item One-hot encoding avec \texttt{np\_utils.to\_categorical}
  \item Compiler (SGD, cross-entropy) et fit
  \item Evaluate sur le test set
\end{enumerate}

Questions pos\'ees~:
\begin{itemize}
  \item Espace des images et sa taille~? ($\mathbb{R}^{784}$)
  \item Nombre de param\`etres du mod\`ele~? (7\,850)
  \item Convexit\'e de la cross-entropy~? (Oui, global)
\end{itemize}

\paragraph{Ex2 -- MLP (2 cellules code)~:}
\begin{enumerate}
  \item \texttt{Dense(100)} + sigmo\"ide
  \item \texttt{Dense(10)} + softmax
  \item Fonction \texttt{saveModel}
\end{enumerate}

Questions~:
\begin{itemize}
  \item Nombre de param\`etres du MLP~? (79\,510)
  \item Convexit\'e du MLP~? (Non, local)
\end{itemize}

\paragraph{Ex3 -- CNN (2 cellules code)~:}
\begin{enumerate}
  \item Reshape en $28\times28\times1$
  \item Architecture LeNet-like (Conv+Pool+Conv+Pool+Flatten+FC)
\end{enumerate}

Questions~:
\begin{itemize}
  \item Temps d'une \'epoque~?
  \item Comparaison CPU vs GPU
\end{itemize}

\subsection{Partie 2 -- Manifold Untangling}

\paragraph{Ex4 -- t-SNE (8 cellules code)~:}
\begin{enumerate}
  \item Imports (matplotlib, sklearn, scipy)
  \item Chargement MNIST
  \item t-SNE sur les donn\'ees de test
  \item D\'efinition \texttt{convexHulls}
  \item D\'efinition \texttt{best\_ellipses}
  \item D\'efinition \texttt{neighboring\_hit}
  \item Calcul des enveloppes/ellipses/NH
  \item D\'efinition \texttt{visualization} et affichage
\end{enumerate}

Question~: les trois m\'etriques sont-elles li\'ees \`a la
s\'eparabilit\'e des classes~? Qu'est-ce qui les diff\`ere~?

\paragraph{Ex5 -- Repr\'esentations internes (3 cellules code)~:}
\begin{enumerate}
  \item D\'efinition \texttt{loadModel}
  \item Chargement du mod\`ele entra\^in\'e
  \item Extraction des repr\'esentations latentes (pop + predict)
    et projection t-SNE
\end{enumerate}

% ===================================================================
\section{Instructions pour l'\'etudiant}
% ===================================================================

\begin{enumerate}
  \item Ex\'ecuter les cellules de code dans l'ordre.
  \item Compl\'eter les cellules marqu\'ees \texttt{\# TODO} ou
    vides.
  \item R\'epondre aux questions dans les cellules markdown.
  \item Sauvegarder les mod\`eles avec \texttt{saveModel}.
  \item Comparer les performances LR (92\,\%) vs MLP (98\,\%)
    vs CNN (99\,\%).
  \item Visualiser le manifold untangling avec t-SNE.
\end{enumerate}

\newpage
\section*{Questions cl\'es (QCM)}
\addcontentsline{toc}{section}{Questions cl\'es (QCM)}

\begin{enumerate}
  \item Quelle librairie est utilis\'ee pour le Deep Learning~?
    \textbf{R. Keras.}
  \item Quel jeu de donn\'ees est utilis\'e~? \textbf{R. MNIST.}
  \item Quelle est la fonction de co\^ut pour la classification~?
    \textbf{R. Cross-entropy.}
  \item Quelle optimisation est utilis\'ee~?
    \textbf{R. SGD (Stochastic Gradient Descent).}
  \item Quelle m\'ethode cr\'ee un r\'eseau s\'equentiel vide~?
    \textbf{R. \texttt{Sequential()}.}
  \item Quelle m\'ethode ajoute une couche~?
    \textbf{R. \texttt{model.add()}.}
  \item Quelle m\'ethode entra\^ine le mod\`ele~?
    \textbf{R. \texttt{model.fit()}.}
  \item Quelle m\'ethode \'evalue le mod\`ele~?
    \textbf{R. \texttt{model.evaluate()}.}
  \item Quelle fonction d'activation pour la classification
    multiclasse~? \textbf{R. Softmax.}
  \item Quelle m\'ethode de visualisation 2D est utilis\'ee~?
    \textbf{R. t-SNE.}
\end{enumerate}

\end{document}
